\section{Which two-factor eta products survive the zero test?}
For positive integers $a\le b$, every $f_{a,b}(z)=\eta(az)\eta(bz)$ satisfies
the strict quadratic Jensen inequalities at all shifts, for coefficients from
$G(w)=\Lambda_{a,b}(1/2+\sqrt w)$, where
$\Lambda_{a,b}(s)=(\sqrt{ab})^s\int_0^\infty f_{a,b}(iy)y^{s-1}\,dy$.
Here $J^{d,n}(X)=\sum_{j=0}^d \binom{d}{j}G^{(n+j)}(0)X^j$.
This completed Mellin transform is zero-free in $\operatorname{Re}s>1$ exactly when
\[
                         b/a\in\{1,2,3,5,7,11,23\}.
\]
At every other ratio the transform has infinitely many zeros there, and
at every fixed shift the number of nonreal Jensen roots tends to infinity
as the degree grows. In the coprime
subfamily with $24\mid a+b$, the sole surviving pair is $(1,23)$.

\paragraph{Proof.}
A common dilation leaves this normalized completion invariant, by a change of variables in the integral.
First divide $a,b$ by their gcd. Now set $g=\gcd(a+b,24)$, $r=24/g$,
$A=ra$, $B=rb$, $N=AB$ and $m=(a+b)/g$. Both eta congruence sums for
$f_{A,B}$ at level $N$ equal $A+B=24m$ (count both factors when $A=B$).
The cusp-order formula has positive summands. Thus this is a weight-one
cusp form of character $\chi(d)=(\frac{-AB}{d})$; see
K.~Ono, \emph{The Web of Modularity}, Theorems 1.64--1.65.
Its Fourier expansion starts with $q^m$ with coefficient one; denote its
coefficients by $a_n$. We have $\gcd(m,N)=1$ since $a+b$ is coprime to
$ab$ and $m$ is coprime to $r$.
If $m>1$, the coefficient of $q$ in $T_m f_{A,B}$ is $a_m=1$ whereas
$a_1=0$. This rules out an eigenfunction for all good-prime $T_p$,
which generate $T_m$ by the Hecke relations.

If $m=1$, $(A,B)$ is one of the twelve two-part partitions of $24$ and
$a_1=1$. The four pairs $(5,19),(7,17),(9,15),(11,13)$ have $a_2=a_4=0$;
the pair $(10,14)$ has $a_3=a_9=0$. The chosen prime $p$ in each case
is coprime to $N$, so the eigenform relation $a_{p^2}=a_p^2-\chi(p)$ fails.
The remaining seven are the classical multiplicative weight-one eta
eigenforms. Their ratios give the set displayed above, and the absolutely
convergent Euler products give nonvanishing on $\operatorname{Re}s>1$.
See Dummit, Kisilevsky and McKay, \emph{Multiplicative products of
$\eta$-functions}, Contemp. Math. 45 (1985), 89--98, and
\href{https://www.numdam.org/item/10.5802/jtnb.495.pdf}{Voskresenskaya, section 2}
for the explicit eigenform list. By
\href{https://research-information.bris.ac.uk/ws/files/95488278/mf_zeros_revised.pdf}
{Booker--Thorne, Theorem 1.1 and Remark (2)}, each excluded form
has infinitely many zeros in $\operatorname{Re}s>1$.

For the Jensen claim, center the Mellin kernel by putting $y=e^u/\sqrt{ab}$.
Apart from a positive constant its logarithm is
$h(u-v)+h(u+v)$, with
$h(t)=t/4+\log\eta(ie^t)$ and $v=\tfrac12\log(b/a)$.
The all-shift quadratic result follows from the paired-eta curvature estimate
of the incoming \texttt{jensen.tex}. For clarity, its derivative estimate uses
the term $-2\pi e^t/24$ and the terms
$-\sigma_{-1}(k)e^{-2\pi k e^t}$ in the expansion of $h$ for $t\ge0$.
For $Y=2\pi k e^t>6$ the fourth derivative of $e^{-Y}$ equals
$e^{-Y}Y(Y^3-6Y^2+7Y-1)>0$.
Eta inversion makes $h$ even, so $h^{(4)}<0$ on the real axis.
For the even logarithm $g$ of the centered kernel this gives
$(g(\sqrt{x}))''<0$ for $x>0$, by integrating $g^{(4)}<0$ first
to obtain $g'''(u)<0$ and then $(ug''-g')'=ug'''<0$ for $u>0$.
The normalized Mellin moments of $e^{g(\sqrt{x})}$ are strictly
log-concave: writing $M_p=\int_0^\infty x^{p-1}e^{-V(x)}dx$ with
$V(x)=-g(\sqrt{x})$, the expectation of $V'$ under the density
$x^p e^{-V(x)}/M_{p+1}$ equals $pM_p/M_{p+1}$ by integration
by parts. This expectation strictly increases with $p$ in steps of one
because weighting a positive density by $x$ increases the mean of the
strictly increasing $V'$ (by strict positive covariance).
The moments $M_{n+1/2}/\Gamma(n+1/2)$ give the Jensen coefficients
up to a constant and a geometric factor.

\paragraph{Failure in every shift.}
Use the normalized modular form $f_{A,B}$, whose completed transform is
$E(s)L(s)$ with $E(s)=(\sqrt{N}/2\pi)^s\Gamma(s)$. The completion above
is invariant under common dilation, by substitution in the Mellin integral.
For every excluded pair choose a closed disk $D\subset\{\operatorname{Re}s>1\}$
with a zero of $L$ inside and none on its boundary. Absolute convergence of
the Dirichlet series gives arbitrarily large $t_j$ with
$L(s+it_j)\to L(s)$ uniformly on $D$. Indeed truncate the series uniformly,
and simultaneously approximate the phases $n^{-it_j}$ by one for the
finitely many indices retained. Every fixed derivative of $L$ is uniformly
bounded on the corresponding vertical strip by absolute convergence slightly
to its left.

Put $w=(s-1/2)^2$ and fix $n\ge0$. Uniformly on $D+it_j$ as $j\to\infty$,
\[
 \frac{[2(s-1/2)]^n G^{(n)}(w)}{E(s)(\log s)^n}
 =L(s)+O_n(1/\log t_j).
\]
To verify this, expand $d/dw=[2(s-1/2)]^{-1}d/ds$ repeatedly.
After the displayed power of $2(s-1/2)$ is applied, the leading term
is $(E L)^{(n)}$; the terms with $k<n$ derivatives have coefficients
$O_n(|s|^{k-n})$. Stirling's differentiated formula gives
$E^{(k)}/E=(\log s)^k+O_k((\log|s|)^{k-1})$ for $k\ge1$ on this strip.
Leibniz's rule proves the estimate. All normalizing factors are nonzero
on the translated disks for large $j$.
Rouch\'e's theorem now gives a zero of $G^{(n)}((s-1/2)^2)$ in every
sufficiently large translated disk. Choose the disks disjoint with increasingly
large positive ordinate: their squares about $1/2$ lie off the real axis and
give infinitely many distinct zeros of $G^{(n)}$.
Finally $J^{d,n}(w/d)\to G^{(n)}(w)$ locally uniformly as $d\to\infty$,
by termwise binomial scaling and domination by the absolute Taylor series.
Any finite collection of disjoint small disks around nonreal zeros must
therefore contain zeros of every sufficiently large scaled polynomial.
This proves the asserted divergence of the nonreal root count.

In fact the number of nonreal roots is $\gg_{a,b,n}\log d/\log\log d$.
To see this, a fixed sufficiently good uniform approximation on $D$ suffices
for Rouch\'e. The corresponding shifts are relatively dense: a finite
Dirichlet sum approximates $L$ uniformly, and the returns of its phase flow
to any neighborhood of the identity have bounded gaps by compactness
of the closure of the flow in its finite torus.
We can therefore choose $\gg T$ disjoint translated disks at heights
between $T/2$ and $T$. On their boundaries the derivative estimate and
Stirling's formula give $|G^{(n)}((s-1/2)^2)|\ge e^{-2T}$ for large $T$.
All the squared disks lie in $|w|\le R$ with $R=O(T^2)$.
Writing the binomial scaling coefficient as
$u_{d,j}=\prod_{k=0}^{j-1}(1-k/d)$ for $j\le d$ and zero otherwise gives
$0\le1-u_{d,j}\le j(j-1)/(2d)$. Positivity of the Taylor coefficients yields
\[
 \max_{|w|\le R}|G^{(n)}(w)-J^{d,n}(w/d)|
 \le R^2G^{(n+2)}(R)/(2d)\le e^{C_nT\log T}/d .
\]
For the last inequality, the centered eta kernel is bounded by
$C\exp(-c e^{|u|})$ from its product and symmetry. Thus
$\log\max_{|w|\le R+1}|G(w)|=O(\sqrt{R+1}\log(R+2))$ by its
Mellin integral (or the gamma integral after substituting $e^{|u|}$).
Cauchy's bound on a unit circle gives the same growth bound for the
fixed derivative. With $T=c_n\log d/\log\log d$ and $c_n$ sufficiently
small, the approximation error is at most $d^{-1/2}<e^{-2T}$.
A second application of Rouch\'e transfers all these disks to the
degree-$d$ polynomial. The constants depend on the eta product and shift.
